\section{The Radon--Nikodym theorem}
\label{sec:radon-nikodym}

The theorem says that absolute continuity is not merely a null-set condition:
whenever $\nu \abscont \mu$, the measure $\nu$ is obtained from $\mu$ by
integrating a density. It is the converse to the elementary fact that
$A \mapsto \int_A f\,d\mu$ defines a measure absolutely continuous with respect
to $\mu$.

\begin{theorem}[Radon--Nikodym]
\label{thm:rn}
Let $\mu$ and $\nu$ be $\sigma$-finite measures on $(\Xset, \salg)$ with
$\nu \abscont \mu$. Then there exists a measurable function
$f : \Xset \to [0, \infty)$ such that
\begin{equation}
  \nu(A) = \int_A f \, d\mu
  \qquad \text{for all } A \in \salg.
  \label{eq:rn}
\end{equation}
The function $f$ is unique up to $\mu$-a.e.\ equality. It is called the
\emph{Radon--Nikodym derivative} of $\nu$ with respect to $\mu$ and written
\begin{equation}
  f = \RN{\nu}{\mu}.
\end{equation}
\end{theorem}

\begin{remark}[Why $\sigma$-finiteness is needed]
Take $\mu$ the counting measure and $\nu$ the Lebesgue measure on $[0,1]$. Then
$\nu \abscont \mu$ (the only $\mu$-null set is $\emptyset$), yet no $f$
satisfies \eqref{eq:rn}: it would need $\int_{\{x\}} f\,d\mu = f(x) = \nu(\{x\})
= 0$ for every $x$, forcing $f \equiv 0$ and hence $\nu \equiv 0$. Counting
measure is not $\sigma$-finite on an uncountable space, and the theorem fails.
\end{remark}

\begin{remark}[The von Neumann proof]
The slick proof avoids constructing $f$ by hand. Assume $\mu, \nu$ finite, set
$\rho = \mu + \nu$, and consider the bounded linear functional
$g \mapsto \int g \, d\nu$ on the Hilbert space $L^2(\rho)$. By Riesz
representation it equals $\int g h \, d\rho$ for some $h \in L^2(\rho)$; algebra
on \eqref{eq:rn}-type identities shows $0 \le h < 1$ $\mu$-a.e.\ (this uses
$\nu \abscont \mu$), and $f = h / (1 - h)$ is the required derivative. The
$\sigma$-finite case follows by exhausting $\Xset$ with finite pieces.
\end{remark}

\subsection{The calculus of Radon--Nikodym derivatives}

The derivative notation is justified by the following rules, each holding
$\mu$- (respectively $\lambda$-) almost everywhere. They are what make
$\RN{\nu}{\mu}$ behave like an ordinary derivative.

\begin{proposition}[Linearity]
\label{prop:rn-linear}
If $\nu_1, \nu_2 \abscont \mu$ are $\sigma$-finite, then
$\nu_1 + \nu_2 \abscont \mu$ and
\begin{equation}
  \RN{(\nu_1 + \nu_2)}{\mu} = \RN{\nu_1}{\mu} + \RN{\nu_2}{\mu}
  \qquad \mu\text{-a.e.}
\end{equation}
\end{proposition}

\begin{proposition}[Change of variables]
\label{prop:rn-integration}
Let $\nu \abscont \mu$ be $\sigma$-finite with density $f = \RN{\nu}{\mu}$. For
every measurable $g : \Xset \to [0, \infty]$,
\begin{equation}
  \int_\Xset g \, d\nu = \int_\Xset g \, f \, d\mu
  = \int_\Xset g \, \RN{\nu}{\mu} \, d\mu.
  \label{eq:rn-cov}
\end{equation}
The identity extends to $\nu$-integrable $g$ of arbitrary sign.
\end{proposition}

\begin{proof}[Proof idea]
Equation \eqref{eq:rn} is exactly \eqref{eq:rn-cov} for indicators
$g = \ind_A$. Extend by linearity to simple functions, then by monotone
convergence to nonnegative measurable $g$, then split $g = g^+ - g^-$ for the
signed case.
\end{proof}

\begin{proposition}[Chain rule]
\label{prop:rn-chain}
If $\nu \abscont \mu \abscont \lambda$ with all three $\sigma$-finite, then
$\nu \abscont \lambda$ and
\begin{equation}
  \RN{\nu}{\lambda} = \RN{\nu}{\mu} \, \RN{\mu}{\lambda}
  \qquad \lambda\text{-a.e.}
  \label{eq:rn-chain}
\end{equation}
In particular, if $\mu \abscont \nu$ and $\nu \abscont \mu$ (equivalent
measures) then $\RN{\nu}{\mu}$ is strictly positive $\mu$-a.e.\ and
\begin{equation}
  \RN{\mu}{\nu} = \left( \RN{\nu}{\mu} \right)^{-1}
  \qquad \mu\text{-a.e.}
\end{equation}
\end{proposition}

\begin{proof}[Proof idea]
Apply Proposition~\ref{prop:rn-integration} with $g = \ind_A \, \RN{\nu}{\mu}$
and the density $\RN{\mu}{\lambda}$:
$\nu(A) = \int_A \RN{\nu}{\mu}\,d\mu = \int_A \RN{\nu}{\mu}\,\RN{\mu}{\lambda}\,
d\lambda$, so the product is a $\lambda$-density for $\nu$; uniqueness in
Theorem~\ref{thm:rn} gives \eqref{eq:rn-chain}. The reciprocal rule is the case
$\lambda = \nu$, where $\RN{\nu}{\nu} = 1$.
\end{proof}
